% Рисунок с доски: функция ползучести материала Кельвина — Фойгта
% Π(t) = (1/2μ)(1 - e^{-t/τ}) с касательной в нуле наклона 1/(2η).
\begin{figure}[htbp]
\centering
\begin{tikzpicture}[>={Latex[length=4pt]}, xscale=1.1]
  \draw[->] (0,0) -- (6.3,0) node[anchor=north east] {$t$};
  \draw[->] (0,0) -- (0,3.1) node[anchor=south east] {$\Pi$};

  \draw[dashed, gray] (0,2.5) -- (6.0,2.5);
  \node[anchor=east, font=\footnotesize] at (-0.05,2.5) {$\dfrac{1}{2\mu}$};

  \draw[very thick, blue!60!black, domain=0:6, samples=120, smooth]
    plot (\x, {2.5*(1-exp(-\x/1.5))});

  % касательная в нуле: наклон Π'(0) = 1/(2η)
  \draw[gray] (0,0) -- (1.5,2.5);
  \draw[dashed, gray] (1.5,0) -- (1.5,{2.5*(1-exp(-1))});
  \node[anchor=north, font=\footnotesize] at (1.5,-0.05) {$\tau$};
  \fill (1.5,{2.5*(1-exp(-1))}) circle (1.2pt);
\end{tikzpicture}
\caption{Функция ползучести $\Pi(t) = \frac{1}{2\mu}(1-e^{-t/\tau})$.
Наклон в нуле $\dt\Pi(0) = \frac{1}{2\mu\tau} = \frac{1}{2\eta}$; за время
$\tau$ набирается $\Pi(\tau) \approx \frac23 \cdot \frac{1}{2\mu}$ — отсюда
название «время запаздывания ползучести».}
\label{fig:creep-function}
\end{figure}
